% arara: xelatex
\documentclass[12pt]{article}

\usepackage{physics}
\usepackage{microtype}
\usepackage{array}
\usepackage{amsmath, amsfonts, amssymb}
\usepackage[top=2cm, left=1.6cm, right=1.6cm, bottom=2cm]{geometry}
\usepackage{booktabs}
\usepackage{enumitem}
\usepackage{fontspec}
\usepackage{polyglossia}
\usepackage{hyperref}

\setmainlanguage{russian}
\setotherlanguages{english}
\setmainfont{Arial}

\hypersetup{colorlinks=true,linkcolor=blue!60!black,urlcolor=blue!60!black}

\DeclareMathOperator{\Cov}{\mathbb{C}ov}
\DeclareMathOperator{\Corr}{\mathbb{C}orr}
\DeclareMathOperator{\Var}{\mathbb{V}ar}
\let\P\relax
\DeclareMathOperator{\P}{\mathbb{P}}
\DeclareMathOperator{\E}{\mathbb{E}}

\newcommand \e{\varepsilon}
\newcommand \cN{\mathcal{N}}
\newcommand{\cF}{\mathcal{F}}
\newcommand{\step}[1]{\medskip\noindent\textbf{#1}\par\nobreak\vspace{2pt}}

\setlength{\parindent}{0pt}
\setlength{\parskip}{6pt}

\title{Midterm 2026 Retake --- решения}
\author{Временные ряды, ВШЭ, весна 2026}
\date{}

\begin{document}
\maketitle

\begin{center}
\textbf{Разбалловка.}\\[4pt]
\begin{tabular}{@{}l|ccc|c@{}}
\toprule
Задача & (a) & (b) & (c) & \textbf{Итого} \\
\midrule
1 --- ETS$(AAN)$                          & 5 & 5 & --- & 10 \\
2 --- GARCH$(1,1)$                       & 5 & 5 & 5   & 15 \\
3 --- Лаговые полиномы и $MA(\infty)$    & 5 & 5 & 5   & 15 \\
4 --- Ковариационная функция             & 10 & 5 & --- & 15 \\
5 --- Нелинейный процесс                 & 5 & 5 & 5   & 15 \\
\midrule
\multicolumn{4}{r|}{\textbf{Всего за работу:}} & \(\mathbf{70}\) \\
\bottomrule
\end{tabular}
\end{center}

\section*{Задача 1. ETS$(AAN)$ \hfill [10 баллов]}

Рассматривается модель
\[
\begin{cases}
u_t \sim \cN(0;\sigma^2),\\
\ell_t = \ell_{t-1} + b_{t-1} + 0.3u_t,\\
b_t = b_{t-1} + 0.2u_t,\\
y_t = \ell_{t-1} + b_{t-1} + u_t.
\end{cases}
\]

\step{(a) ARIMA-представление. \hfill [5 баллов]}
Имеем
\[
y_t=\ell_{t-1}+b_{t-1}+u_t
=\ell_{t-2}+2b_{t-2}+0.5u_{t-1}+u_t.
\]
Так как
\[
y_{t-1}=\ell_{t-2}+b_{t-2}+u_{t-1},
\]
то
\[
y_t=y_{t-1}+b_{t-2}-0.5u_{t-1}+u_t.
\]
Аналогично,
\[
b_{t-2}=y_{t-1}-y_{t-2}+0.7u_{t-2}-u_{t-1}.
\]
Подставляя, получаем
\[
y_t-2y_{t-1}+y_{t-2}=u_t-1.5u_{t-1}+0.7u_{t-2}.
\]
Следовательно,
\[
\boxed{(1-L)^2y_t=(1-1.5L+0.7L^2)u_t.}
\]

\step{(b) Условные моменты. \hfill [5 баллов]}
Из уравнений состояния
\[
y_{t+1}=\ell_t+b_t+u_{t+1}.
\]
При условии на $\cF_t$ величины $\ell_t$ и $b_t$ известны, а $u_{t+1}\sim\cN(0;\sigma^2)$ независима от прошлого. Поэтому
\[
\boxed{\E(y_{t+1}\mid\cF_t)=\ell_t+b_t, \qquad \Var(y_{t+1}\mid\cF_t)=\sigma^2.}
\]

\section*{Задача 2. GARCH$(1,1)$ \hfill [15 баллов]}

Рассматривается модель
\[
r_t=\sigma_t z_t,\qquad z_t\sim\cN(0,1),\qquad
\sigma_t^2=2+0.2r_{t-1}^2+0.7\sigma_{t-1}^2.
\]

\step{(a) Безусловная дисперсия. \hfill [5 баллов]}
Так как
\[
0.2+0.7=0.9<1,
\]
безусловная дисперсия существует. Обозначим
\[
v=\Var(r_t)=\E(r_t^2)=\E(\sigma_t^2).
\]
Тогда
\[
v=2+0.2v+0.7v,
\]
откуда
\[
0.1v=2,\qquad v=20.
\]
Следовательно,
\[
\boxed{\Var(r_t)=20.}
\]

\step{(b) Двухшаговый прогноз волатильности. \hfill [5 баллов]}
При $r_T=3$ и $\sigma_T^2=4$:
\[
\E(\sigma_{T+1}^2\mid\cF_T)=2+0.2\cdot 3^2+0.7\cdot 4=2+1.8+2.8=6.6.
\]
Далее
\[
\E(r_{T+1}^2\mid\cF_T)=\E(\sigma_{T+1}^2 z_{T+1}^2\mid\cF_T)
=\sigma_{T+1}^2 \E(z_{T+1}^2)=6.6.
\]
Поэтому
\[
\E(\sigma_{T+2}^2\mid\cF_T)=2+0.2\cdot 6.6+0.7\cdot 6.6=7.94.
\]
Итак,
\[
\boxed{\E(\sigma_{T+1}^2\mid\cF_T)=6.6,\qquad \E(\sigma_{T+2}^2\mid\cF_T)=7.94.}
\]

\step{(c) Долгосрочный предел. \hfill [5 баллов]}
Обозначим
\[
f_h=\E(r_{T+h}^2\mid\cF_T)=\E(\sigma_{T+h}^2\mid\cF_T).
\]
Тогда для $h\ge 1$
\[
f_{h+1}=2+0.9f_h.
\]
Предел $f_\infty$ удовлетворяет
\[
f_\infty=2+0.9f_\infty,
\]
то есть
\[
f_\infty=20.
\]
Следовательно,
\[
\boxed{\lim_{h\to\infty}\E(r_{T+h}^2\mid\cF_T)=20.}
\]

\section*{Задача 3. Лаговые полиномы и $MA(\infty)$ \hfill [15 баллов]}

Рассматривается уравнение
\[
y_t-0.8y_{t-1}+0.12y_{t-2}=4+u_t-0.2u_{t-1},
\qquad
u_t\sim WN(0,1).
\]

\step{(a) Разложение на множители. \hfill [5 баллов]}
Левая часть:
\[
1-0.8L+0.12L^2=(1-0.6L)(1-0.2L).
\]
Правая часть:
\[
1-0.2L.
\]
Итак,
\[
\boxed{(1-0.8L+0.12L^2)=(1-0.6L)(1-0.2L),\qquad 1-0.2L \text{ уже разложен.}}
\]

\step{(b) Более простое уравнение. \hfill [5 баллов]}
Сначала найдём математическое ожидание стационарного решения:
\[
(1-0.8+0.12)\mu=4 \quad \Longrightarrow \quad 0.32\mu=4 \quad \Longrightarrow \quad \mu=12.5.
\]
Тогда
\[
(1-0.6L)(1-0.2L)(y_t-\mu)=(1-0.2L)u_t.
\]
Для стационарных решений общий множитель $(1-0.2L)$ сокращается, и получаем
\[
(1-0.6L)(y_t-\mu)=u_t.
\]
Эквивалентно,
\[
y_t-0.6y_{t-1}=5+u_t.
\]
Следовательно,
\[
\boxed{(1-0.6L)(y_t-12.5)=u_t}
\quad\text{или}\quad
\boxed{y_t=5+0.6y_{t-1}+u_t.}
\]

\step{(c) Явный вид $MA(\infty)$. \hfill [5 баллов]}
Так как $|0.6|<1$,
\[
y_t-12.5=(1-0.6L)^{-1}u_t=\sum_{j=0}^{\infty}0.6^j u_{t-j}.
\]
Отсюда
\[
\boxed{y_t=12.5+u_t+0.6u_{t-1}+0.6^2u_{t-2}+\cdots.}
\]

\section*{Задача 4. Ковариационная функция \hfill [15 баллов]}

Задано:
\[
\gamma_k^{(1)}=
\begin{cases}
2, & k=0,\\
1, & |k|=1,\\
0, & |k|\ge 2,
\end{cases}
\qquad
\gamma_k^{(2)}=
\begin{cases}
1, & k=0,\\
2, & |k|=1,\\
0, & |k|\ge 2.
\end{cases}
\]

\step{(a) Проверка существования. \hfill [10 баллов]}
Для $\gamma^{(1)}$ рассмотрим процесс
\[
x_t=\eta_t+\eta_{t-1},\qquad \eta_t\sim WN(0,1).
\]
Тогда
\[
\Var(x_t)=1+1=2,\qquad \Cov(x_t,x_{t-1})=1,
\]
а при $|k|\ge 2$ общих шумов уже нет, поэтому
\[
\Cov(x_t,x_{t-k})=0.
\]
Значит, $\gamma^{(1)}$ действительно является ковариационной функцией стационарного процесса.

Для $\gamma^{(2)}$ рассмотрим ковариационную матрицу размера $2$:
\[
\Gamma_2=
\begin{pmatrix}
1 & 2\\
2 & 1
\end{pmatrix}.
\]
Её определитель равен
\[
\det \Gamma_2 = 1\cdot 1 - 2\cdot 2 = -3 < 0.
\]
Ковариационная матрица не может иметь отрицательный определитель, если она должна быть неотрицательно определённой. Поэтому $\gamma^{(2)}$ не может быть ковариационной функцией стационарного процесса.

Итак,
\[
\boxed{\gamma^{(1)}\text{ может быть ковариационной функцией,}\qquad \gamma^{(2)}\text{ не может.}}
\]

\step{(b) Явный пример. \hfill [5 баллов]}
Для допустимой функции подходит процесс
\[
\boxed{x_t=\eta_t+\eta_{t-1},\qquad \eta_t\sim WN(0,1).}
\]
У него
\[
\Var(x_t)=2,\qquad \Cov(x_t,x_{t-1})=1,\qquad \Cov(x_t,x_{t-k})=0 \text{ при } |k|\ge 2,
\]
то есть ковариационная функция совпадает с $\gamma^{(1)}$.

\section*{Задача 5. Нелинейный процесс \hfill [15 баллов]}

Пусть $\e_t$ независимы и
\[
\e_t\sim\cN(0,1),\qquad x_t=\e_t\e_{t-1}.
\]

\step{(a) Слабая стационарность. \hfill [5 баллов]}
Процесс $(x_t)$ является даже строго стационарным, поскольку это одна и та же измеримая функция от сдвигов i.i.d.-последовательности $(\e_t)$.

Проверим моменты:
\[
\E(x_t)=\E(\e_t)\E(\e_{t-1})=0.
\]
Далее,
\[
\Var(x_t)=\E(\e_t^2\e_{t-1}^2)=\E(\e_t^2)\E(\e_{t-1}^2)=1.
\]
Для лага 1:
\[
\Cov(x_t,x_{t-1})=\E(\e_t\e_{t-1}^2\e_{t-2})=0.
\]
Для $|k|\ge 2$ произведение тоже содержит независимый множитель с нулевым средним, поэтому
\[
\Cov(x_t,x_{t-k})=0.
\]
Следовательно,
\[
\boxed{(x_t)\text{ является слабо стационарным процессом.}}
\]

\step{(b) Белый шум. \hfill [5 баллов]}
Мы уже показали, что
\[
\E(x_t)=0,\qquad \Var(x_t)=1,\qquad \Cov(x_t,x_{t-k})=0 \text{ при } k\ne 0.
\]
Значит,
\[
\boxed{(x_t)\text{ является белым шумом.}}
\]

\step{(c) Независимость. \hfill [5 баллов]}
Величины $x_t$ и $x_{t+1}$ делят общий множитель $\e_t$, поэтому независимость неочевидна. Проверим её через квадраты:
\[
x_t^2=\e_t^2\e_{t-1}^2,\qquad x_{t+1}^2=\e_{t+1}^2\e_t^2.
\]
Тогда
\[
\E(x_t^2)=1,\qquad \E(x_{t+1}^2)=1,
\]
а
\[
\E(x_t^2x_{t+1}^2)
=\E(\e_{t-1}^2)\E(\e_t^4)\E(\e_{t+1}^2)
=1\cdot 3\cdot 1=3.
\]
Следовательно,
\[
\Cov(x_t^2,x_{t+1}^2)=3-1\cdot 1=2\ne 0.
\]
Значит, $x_t^2$ и $x_{t+1}^2$ зависимы, а значит и $x_t$, $x_{t+1}$ независимыми быть не могут. Итак,
\[
\boxed{(x_t)\text{ не состоит из независимых случайных величин.}}
\]

\end{document}
